\section{CONFIG-002: Incomplete agent.yaml in Repository}
\label{sec:config-002}

\subsection*{Metadata}
\begin{tabular}{ll}
\textbf{Issue ID} & CONFIG-002 \\
\textbf{Severity} & Medium \\
\textbf{Category} & Configuration \\
\textbf{Branch} & \branchref{fix/CONFIG-002-incomplete-agent-yaml} \\
\textbf{Changes} & \branchdiff{fix/CONFIG-002-incomplete-agent-yaml} \\
\end{tabular}

\subsection*{Executive Summary}
The project's \texttt{agent.yaml} file (used as a template/demo) has an empty \texttt{preferred} model field, making it non-functional out of the box. A user who clones the repository and runs \texttt{gitclaw} without modifying the config sees a confusing error: \texttt{No model configured}. While the scaffolder in \texttt{src/index.ts} creates a correctly populated config, the committed config in the repo should either be functional or contain clear guidance.

The fix addresses both sides: the repository's \texttt{agent.yaml} now defaults to \texttt{openai:gpt-4o-mini}, and the loader validates that \texttt{model.preferred} is not empty, providing a helpful error message when it is.

\subsection*{Root Cause Analysis}

The repository's \texttt{agent.yaml} contained an empty preferred model:

\begin{lstlisting}[language=yaml]
# agent.yaml — Before fix
model:
  preferred: ""
  fallback: []
\end{lstlisting}

The scaffolder code in \texttt{src/index.ts} correctly handles new installations by providing a default model, but the committed file in the repository itself was never updated:

\begin{lstlisting}[language=TypeScript]
// src/index.ts:243-261 — Scaffolder
const defaultModel = model || "openai:gpt-4o-mini";
const yaml = [
    'spec_version: "0.1.0"',
    `name: ${agentName}`,
    "version: 0.1.0",
    `description: Gitclaw agent for ${agentName}`,
    "model:",
    `  preferred: "${defaultModel}"`,
    "  fallback: []",
    // ...
].join("\n");
await writeFile(agentYamlPath, yaml, "utf-8");
\end{lstlisting}

\subsection*{Fix Implementation}

Two changes were made. First, the repository's \texttt{agent.yaml} was given a working default model:

\begin{lstlisting}[language=yaml]
# agent.yaml — After fix
model:
  preferred: "openai:gpt-4o-mini"
  fallback: []
\end{lstlisting}

Second, the loader now validates that \texttt{model.preferred} is not empty, with a descriptive error:

\begin{lstlisting}[language=TypeScript]
// src/loader.ts — After fix: validation
if (!manifest.model?.preferred || manifest.model.preferred.trim() === "") {
    throw new Error(
        "model.preferred is empty or not set.\n" +
        "Set it in agent.yaml:\n" +
        '  model:\n    preferred: "provider:model-id"\n' +
        "Or pass --model provider:model on the command line.",
    );
}
\end{lstlisting}

Additionally, the loader now handles empty files and YAML parse errors gracefully:

\begin{lstlisting}[language=TypeScript]
// src/loader.ts — After fix: parse validation
const parsed = yaml.load(manifestRaw);
if (!parsed || typeof parsed !== "object") {
    throw new Error(
        "agent.yaml does not contain a valid configuration object.\n" +
        "Ensure the file starts with a top-level mapping (e.g., 'name: my-agent').",
    );
}
\end{lstlisting}

\subsection*{Affected Files}
\begin{itemize}
    \item \srcfile{agent.yaml} --- fixed empty preferred model
    \item \srcfile{src/loader.ts} --- added model validation and parse error handling
    \item \srcfile{src/\_\_tests\_\_/config-002.test.ts} --- new validation tests
\end{itemize}

\subsection*{Verification}
\begin{lstlisting}[language=TypeScript]
// config-002.test.ts
test("agent.yaml has a valid preferred model", () => {
    const repoConfig = yaml.load(readFileSync("agent.yaml", "utf-8")) as any;
    ok(repoConfig.model?.preferred, "agent.yaml should have a preferred model");
    ok(repoConfig.model.preferred !== "", "model.preferred should not be empty");
    ok(repoConfig.model.preferred.includes(":"),
        "model.preferred should be in provider:model format");
});

test("loadAgent validates empty manifest", async () => {
    const { loadAgent } = await import("../loader.ts");
    try {
        await loadAgent("/nonexistent", undefined);
        ok(false, "should have thrown");
    } catch (err: any) {
        ok(err.message.includes("agent.yaml") ||
            err.message.includes("model.preferred"));
    }
});
\end{lstlisting}
